\subsection{Definición del problema}

El \textit{Capacitated Vehicle Routing Problem} (CVRP) consiste en diseñar un conjunto de
rutas de costo mínimo para abastecer a un conjunto de clientes con demanda conocida,
empleando una flota homogénea de vehículos con capacidad limitada que parten y retornan a
un depósito central. Cada cliente debe ser visitado exactamente una vez por un único
vehículo, y la demanda acumulada de los clientes asignados a una ruta no puede exceder la
capacidad del vehículo.

\subsection{Formulación matemática}

Sea $G=(V,E)$ un grafo completo no dirigido, donde $V=\{1,2,\dots,n\}$ es el conjunto de
nodos, con el nodo $1$ representando el depósito, y $E$ el conjunto de aristas. Cada
cliente $i\in V\setminus\{1\}$ posee una demanda $d_i \ge 0$, y cada arista
$\{i,j\}\in E$ tiene un costo $c_{ij}=c_{ji}$ (problema simétrico). Se considera una flota
homogénea de $m$ vehículos, todos con capacidad $D$. Una solución es un conjunto de rutas
$\mathcal{R}$, cada una un ciclo que parte y retorna al depósito, tal que cada cliente es
visitado exactamente una vez y la demanda acumulada de cada ruta no excede $D$.

La función objetivo minimiza la distancia total recorrida por la flota:
\begin{equation}
    \minimo \; z = \sum_{r \in \mathcal{R}} \sum_{\{i,j\} \in r} c_{ij},
    \label{eq:m2_objetivo}
\end{equation}
sujeta a las condiciones de factibilidad: cada cliente pertenece exactamente a una ruta
(las rutas particionan $V\setminus\{1\}$) y cada ruta $r\in\mathcal{R}$ respeta la
capacidad,
\begin{equation}
    \sum_{i \in r} d_i \le D, \quad \forall\, r \in \mathcal{R}.
    \label{eq:m2_capacidad}
\end{equation}

Esta es la misma definición del CVRP abordada en el Informe~1. La diferencia radica en el
método de resolución: mientras el enfoque exacto formulaba el problema como un
\textit{problema lineal entero} y lo resolvía sobre un grafo no dirigido, en este trabajo
la metaheurística PSO no resuelve el modelo entero de forma directa, sino que explora el
espacio de soluciones evaluando \eqref{eq:m2_objetivo} como función de aptitud y
garantizando \eqref{eq:m2_capacidad} en la etapa de decodificación (Capítulo~\ref{cap:metodo}).
